\documentclass[aps,pra,twocolumn,superscriptaddress,longbibliography,floatfix,nofootinbib]{revtex4-2}

\usepackage{amsmath,amssymb,amsfonts}
\usepackage{graphicx}
\usepackage[dvipsnames]{xcolor}
\usepackage{hyperref}
\usepackage{booktabs}
\usepackage{bm}
\usepackage{mathtools}

\providecommand{\ket}[1]{\lvert #1 \rangle}
\providecommand{\bra}[1]{\langle #1 \rvert}
\newcommand{\Id}{\openone}
\newcommand{\Fnorm}[1]{\lVert #1 \rVert_{F}}
\newcommand{\Opnorm}[1]{\lVert #1 \rVert_{\mathrm{op}}}
\newcommand{\Maxnorm}[1]{\lVert #1 \rVert_{\max}}

\hypersetup{
  colorlinks=true,
  linkcolor=blue!60!black,
  citecolor=blue!60!black,
  urlcolor=blue!60!black,
  pdfauthor={Berkay Y{\"u}ksel Sayim},
  pdftitle={Resolving the Full Modular Data of a Doubled-Fibonacci
            String-Net on a Finite Torus by Point-Group Rotation, and an
            Obstruction to Quasi-One-Dimensional Anyon Exchange},
  pdfsubject={cond-mat.str-el},
  pdfkeywords={string-net, Levin-Wen, doubled Fibonacci, modular data,
               topological order, anyons, exact diagonalization,
               superselection, holonomy}
}

\begin{document}

\setcounter{secnumdepth}{3}

\title{Resolving the Full Modular Data of a Doubled-Fibonacci String-Net
       on a Finite Torus by Point-Group Rotation, and an Obstruction to
       Quasi-One-Dimensional Anyon Exchange}

\author{Berkay Y\"uksel Sayim}
\email{berksa@tutamail.com}
\affiliation{Independent Research, Germany}
\affiliation{\href{https://orcid.org/0009-0004-4993-7352}{ORCID:
0009-0004-4993-7352}}

\date{2026-09-22}

\begin{abstract}
We report a real-space resolution of the full modular data of a doubled-Fibonacci Levin--Wen
string-net, obtained by exact diagonalization on finite torus supercells ($2\times2$,
$2\times3$, $3\times3$), together with a sharply scoped negative result for the dynamical half
of the same program --- and we keep the classes of evidence strictly separate. (i)~The vacuum row of the modular
$S$-matrix, $|S_{0a}| = d_a/D$, is confirmed by \emph{two structurally distinct
lattice routes} --- an operator route built from a flux basis, and a
Verlinde--Casimir route that reads the quantum dimensions from the lattice loop
spectrum without using the $F$ or $R$ symbols --- the latter reproducing the
analytic values at $1.1\times10^{-16}$ and the two agreeing at $<10^{-9}$, the
precision at which the operator-route values are recorded. (ii)~The ground-state
degeneracy $\mathrm{GSD}=4$ is a separate check, invariant across $2\times2$,
$2\times3$ and $3\times3$; $|S|$ is invariant under $24$ sampled $O(4)$ gauge
transformations of the ground-state manifold and size-robust to
$\le2.5\times10^{-11}$. (iii)~The Wilson-loop algebra \emph{alone} does not fix
the chiral sign, and we \emph{measure} that limit rather than assert it: the
loop algebra is $\mathbb{C}\oplus M_3(\mathbb{C})$ ($\dim=10=1^{2}+3^{2}$), its
antisymmetric part vanishes identically on the real-symmetric holonomy, and the
chiral splitting capacity is zero. This is a concrete instance of a theorem of
Li and Mong, who also name the remedy: a point-group rotation of the
ground-state manifold, real but \emph{not} symmetric --- precisely the property the loop route
lacks --- available on every $L\times L$ supercell. Carrying it out, we extract
the full signed $S$ and $T$ from the lattice on \emph{two} sizes, with no
analytic doubled-Fibonacci $S$ or $T$ supplied to the solver, up to the
$\mathbb{Z}_2$ relabeling intrinsic to an achiral theory; a validated
two-channel error budget with no free parameters accounts on both sizes for
the reference error --- the distance to that analytic target, formed in a
separate scoring step, not an output of the solver.
The modular data themselves are not new with this
work --- Francuz and Dziarmaga obtained them for this model in 2020; new here
are the finite-lattice demonstration of the rotation channel, the quantified
error budget, and the measured loop-algebra limit as a concrete instance of the
Li--Mong bound. Complementing this static certificate, we show that on the
one-dimensionally mapped anyonic ladder the Hamiltonian transport of a single
mobile Ising anyon around a pinned object cannot realize the non-Abelian braid
matrix: the closed-loop holonomy is exactly the identity and
$[U(C_1),U(C_2)]=0$, established through two independent obstructions, three
null escape routes, and a machinery-alive control. That obstruction is
structural and independent of both the anyon type and the system size, so it
constrains the Fibonacci ambition motivating this work and not merely the Ising
tracer that exhibits it. Both halves were produced under an adversarial
verification harness that, in the course of this work, retracted an attribution of its own ---
crediting the lattice with the signed data before a channel existed to carry it --- and then
restored it with proof through the rotation channel. The values were never in question, only
their provenance.
\end{abstract}

\maketitle

\section{Introduction and framework}
\label{sec:intro}

\subsection{The original goal}
\label{sec:goal}

This work began in pursuit of an ambitious target: a real-space braiding
simulator. Concretely, to carry two non-Abelian Fibonacci anyons ($\tau$)
around one another \emph{in real space} and to measure the non-Abelian
Berry/monodromy matrix they accumulate, whose target value is
\begin{equation}
  M_{\tau\tau} = -1/\varphi^{2} = -0.382\,,
  \label{eq:target}
\end{equation}
with $\varphi=(1+\sqrt5)/2$ the golden ratio. This quantity is a characteristic
monodromy diagnostic of Fibonacci statistics; universality rests on the
matrix-valued braid action, of which this scalar is one invariant. Obtaining it
\emph{Hamiltonian-dynamically},
through genuine quasiparticle transport rather than as a digital gate sequence
on a qubit chip, would be categorically distinct from the recent digital
demonstrations of Fibonacci anyons on quantum processors~\cite{Xu2024}, which
realize the statistics as gate circuits with no static Hamiltonian and no
transport dynamics.

The goal splits cleanly into a static and a dynamic half, and the line between
them is where this paper sits. The static half is delivered: we extract, from a
genuine finite-lattice computation, the basis-invariant modular content of the
doubled-Fibonacci Levin--Wen string-net~\cite{LevinWen2005}. The dynamic half
is a verified no-go on this model: the transport is structurally trivial. Both
outcomes were anticipated as publishable before the decisive test was run,
which is why we report them together.

\subsection{The narrow ridge, told plainly}
\label{sec:ridge}

The approach to the dynamic half was sound, and stating that plainly is part of
the contribution. Fibonacci dynamics live in a small fusion space that is
tractable locally via anyonic matrix product states~\cite{Singh2014}, at a cost
scaling as $\chi^{3}$ rather than exponentially, so no large-scale computing
resources were required. A tracer-first ladder of validations was passed
against known answers, and both target braid matrices --- the Fibonacci
monodromy of Eq.~\eqref{eq:target} and the Ising tracer
$e^{-i\pi/4}\sigma_{x}$ --- were built exactly from the framework's own $F$ and
$R$ symbols.

That left one decisive question: does the transport actually realize the
braiding? The clean test was an Ising-Berry gate --- transport a free-fermion
Ising anyon, whose braid answer is known in closed form, and check whether the
measured matrix hits it. Either branch was publishable: a hit would carry the
method through to Fibonacci, a miss would fix the boundary honestly. The single
real hazard was asymmetric: the hopping Hamiltonian is given only
diagrammatically in the source papers~\cite{Singh2014,Ayeni2016}, so it had to
be reconstructed from the originals rather than guessed, because a mis-built
anyon returns a plausible-looking but meaningless number.

The decisive test resolved to the no-go of Sec.~\ref{sec:nogo}.

\subsection{Why an Ising test settles a Fibonacci question}
\label{sec:bridge}

Section~\ref{sec:goal} motivates Fibonacci; Sec.~\ref{sec:nogo} tests Ising. The
bridge is that \emph{the obstruction is structural, and independent of the
anyon type and of the system size}. The two mechanisms established in
Sec.~\ref{sec:nogo} are gauge-removability of the braid $R$-phase on a simply
connected graph, and superselection of the total fusion charge. Neither invokes
the specific fusion category: the first is a statement about closed walks on a
graph, the second about the block structure of every local operator in the
model. Neither depends on the chain length $N$; we verify the flux statement
explicitly on $N=3,5,7$.

Ising is therefore not a substitute target but the \emph{decisive} test: its
braid answer is known in closed form, so a null result cannot be attributed to
an unknown target or to a mis-tuned parameter. The no-go consequently
constrains the Fibonacci ambition that motivates this work, and not merely the
Ising tracer that exhibits it.

\subsection{The harness}
\label{sec:harness}

Both halves were produced under an adversarial verification harness: every
positive claim must survive a battery of locks, each lock must carry a positive
control that makes it fail on a known-nontrivial input, and any claim that
cannot be reproduced from the deposited artifacts is withdrawn. The harness is
not decoration. In the course of this work it corrected an attribution of its
own, twice. An earlier internal headline had credited the lattice with the
\emph{full, signed} modular $S$ without proof; an independent recomputation
withdrew that credit, showing that the raw signed matrix of the operator route
is not symmetric, is not gauge-equivalent to the analytic $S$, and fails the
modular consistency condition at $(ST)^{3}-S^{2} = 1.27$ --- facts that remain
true of that route. A follow-up then over-corrected, declaring the lattice
unable to carry the sign at all; a point-group rotation of the ground-state
manifold, until then untested, restored the attribution with proof.
Section~\ref{sec:cut} traces both corrections and the channel that resolved
them. What moved was never the values, only their provenance --- a discipline
we regard as the honest core of the paper rather than a damage report.

\subsection{Model}
\label{sec:model}

We work with the doubled-Fibonacci Levin--Wen string-net~\cite{LevinWen2005} on
an $L\times L$ torus. The plaquette operators $B_p$ are built from the (real)
Fibonacci $F$-symbols; the ground space is the joint $+1$ eigenspace of the
$B_p$. The anyon content is the Drinfeld center, four sectors
$a\in\{(\Id,\Id),(\tau,\Id),(\Id,\bar\tau),(\tau,\bar\tau)\}$ with quantum
dimensions $\{1,\varphi,\varphi,\varphi^{2}\}$ and total quantum dimension
$D=\sqrt{\sum_a d_a^{2}}=1+\varphi^{2}=3.618\ldots$, so that
$D^{2}=13.09\ldots$. Exact diagonalization is performed at $2\times2$,
$2\times3$ and $3\times3$, of Hilbert-space dimension $175$, $2250$ and
$106250$ respectively.

We stress a scope point up front, because it bounds everything below: a second
lattice \emph{size} beyond $3\times3$ is not reachable by exact
diagonalization. A $4\times4$ state requires ${\sim}13$~GB per vector and
$5\times5$ about $1.4$~PB per vector. Where a claim rests on $3\times3$ alone
we say so.

\section{Lattice-emergent modular content}
\label{sec:c1}

The central discipline of this section is that the classes of evidence are
\emph{not} interchangeable, and we label each one: (a) the vacuum row, fixed by
two structurally distinct lattice routes; (b) the topological spins, which the
loop overlaps constrain only algebraically; and (c) the ground-state
degeneracy, a separate lock. The full signed data rest on a fourth class, the
point-group rotation channel, which we develop in Sec.~\ref{sec:cut}. Bundling
these into a single ``independent routes confirm our result'' would be an
overclaim; keeping them apart is what let the harness catch one earlier in this
work.

\subsection{Class (a): the vacuum row, by two distinct lattice routes}
\label{sec:classa}

Only for the vacuum row of $S$ do two structurally distinct \emph{lattice}
routes exist, and they agree.

The \emph{operator route} builds the commutant of $\{B_p\}$, forms the
$K$-even and $K$-odd ground-space actions of a Wilson-loop operator $W$ along a
noncontractible cycle, diagonalizes them jointly to obtain a flux basis, and
reads $S=U_i^{\dagger}U_j$ as the overlap of the flux bases of the two cycles.

The \emph{Verlinde--Casimir route} shares none of those code paths. It resolves
a four-sector basis from a single hermitian lattice Casimir
$H_W = 0.7\,M_i - 0.3\,M_j$, where $M_g=V^{\mathsf T}W_gV$ is the ground-space
action of the cycle loop; because the two cycle loops fail to commute at finite
size, $\Maxnorm{[M_i,M_j]} = 5.6\times10^{-2}$ (largest absolute entry),
this noncommutativity lifts
the per-cycle degeneracy. It then reads $S_{0a}=d_a/D$ directly from the
lattice loop spectrum via the Verlinde relation~\cite{Verlinde1988}, rather
than from any basis overlap. Its extraction inputs are only
$\{V_{\mathrm{gs}}, W_{\text{cycle }i}, W_{\text{cycle }j}\}$; no $F$, no $R$,
and no analytic $S$ enters.

We separate the two comparisons, because they are recorded at different
precision. The Verlinde--Casimir route, independent of analytic $F/R$ symbols
and of the target modular matrix, reproduces the analytic
values $\{0.2764,0.4472,0.4472,0.7236\}$ at $1.1\times10^{-16}$. The two
\emph{routes} agree with each other at $<10^{-9}$ --- that bound, and not the
$10^{-16}$ figure, is what the operator-route data support, since those values
are recorded to four digits. Conflating the two would overstate the
route-to-route agreement by seven orders of magnitude. This is the only claim
of two structurally distinct lattice routes made anywhere in this paper, and it
applies to the vacuum row alone.

\subsection{Class (b): topological spins --- algebraic from the loops,
lattice-determined by rotation}
\label{sec:classb}

The topological spins are
$T=\mathrm{diag}(1,e^{4\pi i/5},e^{-4\pi i/5},1)$, i.e.\ spins
$\{0,0,\pm4\pi/5\}$. From the loop overlaps \emph{alone} they are fixed only
noncircularly and algebraically --- no analytic $T$ is inserted --- by two
consistency paths applied to the lattice $S$; the direct lattice
\emph{measurement} of $T$ comes not from the loops but from the point-group
rotation channel of Sec.~\ref{sec:cut}. The two consistency paths are:

\begin{enumerate}
\item[(b1)] a \emph{Verlinde fixed point}: solving $(ST)^{3}=S^{2}$ from random
grid starts (residual $3.98\times10^{-11}$). This is a consistency check, and
by itself it does not pin $T$ uniquely, since $S$ leaves Galois orbits open.
\item[(b2)] a \emph{Gauss-sum/anomaly} argument~\cite{Vafa1988,Bantay1997,Kitaev2006}:
$p_{+}=\sum_a d_a^{2}\theta_a = D$ at $c=0$ pins the solution uniquely without
invoking $(ST)^{3}$, and rejects the order-3 alternative ($5.236 \ne D$). The
Galois-conjugate solution is excluded not by this argument but by the
positivity of the extracted quantum dimensions ($d_{\tau}=\varphi>0$ rather
than $-1/\varphi$), i.e.\ by unitarity~\cite{Freedman2012}: Galois conjugation
preserves the order of $T$, so a congruence of level $5$ cannot separate
members of one Galois orbit~\cite{NgSchauenburg2010,DongLinNg2015}.
\end{enumerate}

Both paths take the lattice $S$ as input and reason about it algebraically. Two
facts about the loops explain why this is as far as they reach, and they set up
the channel that reaches further. Real-symmetric lattice loops are phase-free,
so a Wilson-loop overlap cannot carry a spin; and the geometric Dehn twist
available on a $3\times3$ torus has order~3, whereas the theory's $T$ has
order~5, so that twist cannot generate $T$ either --- a directly geometric
order-5 twist would require $L\equiv0\bmod5$. Neither obstacle touches the
point-group rotation $ST$, a \emph{different} modular element whose order is $3$
intrinsically and on every size; the order-3 limitation that blocks the Dehn
twist does not apply to it. That rotation is the channel of Sec.~\ref{sec:cut}.

\subsection{Class (c): ground-state degeneracy}
\label{sec:classc}

$\mathrm{GSD}=4$ is a separate check and a separate class of evidence. The
degeneracy and the gap are invariant across all three lattice sizes: four
ground states, clean gap $1.0$.

The ground-state \emph{energy} is not, and we state the scope rather than let a
reader infer one. $E_0$ is \emph{extensive}: it counts one unit per plaquette,
so $E_0=-4$ on $2\times2$ and $E_0=-9$ on $3\times3$. It is therefore the
degeneracy and the gap --- not the energy --- that carry the invariance, and
quoting a single $E_0$ for the model would be meaningless.

\subsection{Soundness of the $|S|$ extraction}
\label{sec:soundness}

Two properties support the $|S|$ extraction. They are properties of $|S|$; they
do not extend to classes (b) or (c), each of which stands on its own evidence.

\emph{Gauge invariance.} $|S|$ is invariant under $V\to V\cdot Q$ over $24$
sampled $O(4)$ gauge transformations of the ground-state manifold, with a worst
deviation of $2.4\times10^{-15}$ on $2\times2$ and $5.6\times10^{-15}$ on
$2\times3$. That $O(4)$ is the correct residual gauge group is itself a
lattice fact: the $F$-symbols are real, hence the Hamiltonian is real and
time-reversal invariant, hence the $K$-odd part vanishes ($6.9\times10^{-17}$)
and the ground states may be chosen real. The corresponding $U(4)$ test
\emph{fails} at $0.603$ and thereby acts as a discriminator: a test that passes
for every group would be measuring nothing.

\emph{Size robustness.} $|S|$ agrees across $2\times2$, $2\times3$ and
$3\times3$ to $\le2.5\times10^{-11}$, and the $|S|$ multiset
$\{0.2764,0.4472,0.7236\}$ reproduces the doubled-Fibonacci values.

The string operator underlying the extraction is genuine in a checked sense:
it commutes with every plaquette operator ($\le1.7\times10^{-16}$), it has
support outside the ground space, and it acts nontrivially on the ground-state
manifold. Its diagonal Wilson line reproduces the doubled-Fibonacci eigenvalues
$\{\varphi^{2},-1,-1,1/\varphi^{2}\}$ exactly, with error $0.0$. The entry
$+1/\varphi^{2}$ is a static lattice quantity; it must not be confused with the
dynamical target $M_{\tau\tau}=-1/\varphi^{2}$ of Eq.~\eqref{eq:target}, which
carries the opposite sign, belongs to the other half of the program, and was
never measured (Sec.~\ref{sec:nogo}). We call the operator
\emph{chirality-discriminating} (equivalently, $R$-phase-sensitive) rather than
chiral: the theory is achiral, $c=0$.

Honesty requires one further statement here. The $|S|$/quantum-dimension
extraction itself is \emph{single-route on $3\times3$}: it carries no
two-route redundancy at that size, and $5\times5$ is infeasible by exact
diagonalization. What supports it is the genuine out-of-ground-state,
chirality-discriminating string, not a second independent measurement.

\subsection{Consistency locks}
\label{sec:locks}

The categorical data underlying $B_p$ satisfy the hexagon equation to
$2.4\times10^{-16}$ and the pentagon equation with zero violations, and the
Levin--Wen hermiticity, fusion, projector and commutation properties hold at
${\sim}10^{-16}$. These locks are not tautological: a deliberately naive
evaluation of the same expression, $FRF-RFR$, returns $1.18$ rather than zero,
so the lock discriminates.

\section{The honest cut: what the loops cannot give, and how the lattice
gives it}
\label{sec:cut}

\subsection{The loop route, measured: an instance of the Li--Mong bound}
\label{sec:wedderburn}

The operator route of Sec.~\ref{sec:classa} fixes the magnitudes $|S|$ but not
the chiral sign, and we measure that limit rather than assert it. Two facts
establish it.

First, the raw signed matrix that route produces is not a valid modular $S$: it
is not symmetric (symmetry error $1.46$, whereas a valid $S$ must be
symmetric), it is not gauge-equivalent to the analytic $S$ (best over $24$
permutations and $16$ sign choices: $1.08$, against a requirement of
${\sim}10^{-2}$), it fails modular consistency at $(ST)^{3}-S^{2}=1.27$, and the
Verlinde relation applied to it returns negative, noninteger fusion
multiplicities; the corresponding canonical-$S$ lock returns \textsc{false}.
These are honest results of the project's own audit, and they remain true of
that route. A selector-backed variant of it did reach a $1.16\times10^{-10}$
match to the analytic $S/T$ --- but only by \emph{supplying} the sign from an
$SL(2,\mathbb{Z})$ solvability argument rather than measuring it; we report that
number only as a property of the superseded route.

Second, the reason is structural, and we can exhibit it. On $3\times3$ the cycle
holonomies $M_i,M_j$ are exactly real-symmetric, and their antisymmetric parts
vanish identically ($\Fnorm{i(M-M^{\mathsf T})/2} \approx 1.2\times10^{-16}$
and $7.5\times10^{-17}$). The loop algebra $\langle M_i,M_j\rangle$ has
dimension $10=1^{2}+3^{2}$, i.e.\ $\mathbb{C}\oplus M_3(\mathbb{C})$: with only
$\{\Id,\tau\times\bar\tau\}$ Wilson loops available, the four anyon sectors
cannot be separated into four one-dimensional blocks. Each real-symmetric loop
resolves only three distinct levels, the pair $(\tau,\Id)/(\Id,\bar\tau)$
remaining degenerate, and the chiral splitting capacity of that pair is $0.0$. A
single hermitian Casimir does resolve the four-sector \emph{basis}, but a basis
is not a signed $S$.

This is not a shortfall of our implementation; it is a concrete instance of a
theorem. Li and Mong~\cite{LiMong2022} prove that Wilson-loop and
minimum-entropy-state overlaps determine only the absolute-value structure of
the modular data, leaving the individual spins $\theta_a$ inaccessible. Our
measured loop algebra realizes that bound on a finite lattice. The remedy is
theirs too, and the rest of this section takes it up.

One internal quantity must be named to prevent a spurious contradiction. A
separate ``$e{-}15$ match'' of $|S|$ row multisets used \emph{different} row and
column permutations; such a match constrains the multiset of absolute values
only, is not a representation-preserving anyon relabeling, and is a different
object from the $1.16\times10^{-10}$ above. The two must not be conflated ---
this ``value correct, scope wrong'' distinction is one the harness enforces
throughout.

\subsection{The rotation channel}
\label{sec:rotation}

The remedy Li and Mong name is a point-group rotation of the ground-state
manifold, available on every $L\times L$ supercell. On our lattice it is the
$120^{\circ}$ rotation $R$ about a plaquette center; realized on the ground
space as $M_R=V_{\mathrm{gs}}^{\mathsf T}RV_{\mathrm{gs}}$, it is a genuine
symmetry --- orthogonal to $2.96\times10^{-15}$ on $3\times3$ and
$4.41\times10^{-15}$ on $2\times2$, of order three --- and it does \emph{not}
lie in the loop algebra: adjoining it raises the algebra dimension from $10$ to
$16$.

Why the loops missed it is the most instructive line of this paper. A
real-symmetric operator has real eigenvalues and cannot carry a phase; the
Wilson-loop holonomies are exactly such operators, which is why they fix only
$|S|$. A rotation is real but \emph{not} symmetric, so it can. The earlier
statement that ``the $3\times3$ geometry is real-symmetric and therefore
phase-free'' was true of the Wilson-loop observable, not of the lattice --- the
lattice carries other real operators that are not symmetric, and $M_R$ is one.

Because $M_R$ commutes with neither $M_i$ nor $M_j$, it lifts the degeneracy
the loops leave. From the measured pair $(M_{\mathrm{Wilson}},M_R)$ we solve for
the modular data, the rotation supplying the phase and the loop spectrum the
magnitudes. On $3\times3$ the solver returns four solutions per computed
rotation center, with residuals $2.9\times10^{-10}$ to $4.8\times10^{-10}$ and
the correct sign structure on every one, yielding $T=\{e^{\mp4\pi i/5},1,1\}$.

Those centers are not independent repetitions, and we say so rather than let
the count suggest otherwise. On an $L\times L$ torus the nine possible centers
generate only three distinct edge permutations; centers sharing a permutation
perform the same computation and their projected rotations therefore agree bit
for bit, while across all nine the projected rotation agrees to
$1.4\times10^{-15}$, the floating-point floor. Two of the three centers used
here fall in the same class, so the twelve solutions are eight distinct ones.
We deposit one rotation matrix per class and a script that checks this from the
record. Physically the content is that lattice translations act trivially on
the ground space, as topological order on the torus requires; a
center-dependent modular datum would have been the surprise. Running further
centers is therefore a consistency check on the optimizer, not an independent
repetition --- the independent repetition in this work is the second lattice
size.

On $2\times2$ the solver returns $9$ solutions across three centers, residuals
$6.4\times10^{-10}$ to $1.8\times10^{-9}$, again with correct sign structure; the extracted $T$ spectra agree across all
three centers at the level of the extraction residuals themselves, up to the
intrinsic $\mathbb{Z}_2$ relabeling ($\tau\leftrightarrow\bar\tau$) --- verifiable
from \texttt{p6\_extraction\_2x2.json}. The full signed $S$ and $T$ are thereby resolved from
the lattice on two sizes. These are two \emph{sizes} of one lattice route, not
two routes: the signed data rest on the single rotation channel, and no second
structurally distinct lattice route is claimed for them.

The extractor merits a precise description, since it is released with the
paper. It solves a nonlinear least-squares problem whose residual encodes
generic modular-tensor-category axioms --- $S$ symmetric, a real nonnegative
vacuum row, $|\theta_a|=1$, fusion-character consistency, $S^{2}$ equal to
charge conjugation, and $(ST)^{3}=\zeta S^{2}$ with $|\zeta|=1$ (the
central-charge phase, left free). These axioms \emph{define} the object sought
and are applied identically to every input. No analytic $S$ or $T$, no $F$ or
$R$ symbol, and no solvability selector is fed in; a poisoning test confirms the
theory helpers are never reached. Verlinde integrality and $S$-unitarity are
applied as output gates after the solve, and an independent referee harness
rechecks the result against separately held thresholds. The chiral sign enters
through the measured rotation $M_R$; the modular axioms only sieve out
nonmodular branches. This is the structural difference from the superseded
route, where the sign was \emph{supplied} by a solvability selector rather than
\emph{measured}. That the rotation must equal $V(S\cdot T^{*})V^{\dagger}$ up to
a permutation is itself a consequence of the Li--Mong analysis, not an
assumption: the only rotation action consistent with the fusion data is a
permutation preserving both quantum dimension and spin, with Wilson phases
drawn from the Abelian subcategory; for doubled Fibonacci that subcategory is
trivial and the permutation is forced to the identity, so this form is the
unique one allowed.

\subsection{Why the extraction is unique for doubled Fibonacci}
\label{sec:unique}

That the extraction returns a unique answer for doubled Fibonacci is luck
rather than merit --- and it is again Li and Mong's theorem, not our inference,
that says so. They prove that the number of fusion phases left free by the
loop-plus-rotation data equals the number of \emph{self-dual} Abelian anyons.
Doubled Fibonacci has quantum dimensions $\{1,\varphi,\varphi,\varphi^{2}\}$, so
its only Abelian anyon is the vacuum---which is self-dual---giving a count of
$1$: the single fusion phase is trivial, and $S$ and $T$ are fixed. Fibonacci is simply the most favorable
case of their theorem.

The freedom is real when $|\mathcal{A}|>1$. A related but distinct degeneracy
shows up empirically in our own inverse problem: on the fully Abelian toric
code the extractor returns \emph{two} solution classes rather than one
(residuals $1.0\times10^{-15}$ and $1.3\times10^{-15}$), while on doubled
Fibonacci it returns a single class. This nonuniqueness is a degeneracy of
\emph{our} data (one loop spectrum plus $M_R$) rather than an instance of Li
and Mong's fusion-phase freedom, and it arises because those data do not
separate the two classes; its absence on Fibonacci is the strongest evidence
that the doubled-Fibonacci extraction is well posed.

\subsection{Which category the data selects}
\label{sec:whichcategory}

Among the unitary modular tensor categories of rank 4 classified in
Ref.~\cite{RSW2009}, exactly one is compatible with the extracted data:
$Z(\mathrm{Fib})=\mathrm{Fib}\boxtimes\overline{\mathrm{Fib}}$. The quantum
dimensions alone do not suffice --- $\mathrm{Fib}\boxtimes\mathrm{Fib}$ and
$\overline{\mathrm{Fib}}\boxtimes\overline{\mathrm{Fib}}$ share them exactly
--- and are separated only by the topological spins, whose multiset
$\{0,0,\pm4\pi/5\}$ is reproduced by the achiral combination alone. The
residual $\mathbb{Z}_2$ is an automorphism of that single category (exchange of
the two chiral factors), not a competing candidate. All excluded competitors
differ by $O(1)$, so the conclusion is independent of the error budget.

\subsection{The error budget}
\label{sec:errorbudget}

The quantity budgeted here is the reference error: the distance between the
extracted data and the analytic doubled-Fibonacci target, formed in a
separate scoring step and not an output of the solver. It is not a single
number but a two-channel budget, and validating it across both lattice sizes
with no free parameters is a result in its own right. This residual follows
\begin{equation}
  \text{residual} \approx 2.081\,\varepsilon_{\mathrm{frame}} + k\,s^{2},
  \label{eq:budget}
\end{equation}
where $s=\varepsilon_{\mathrm{ideal}}/(3.2668\times10^{-3})$ is the loop
ideality relative to its $3\times3$ value, with two independently measured
channels. The eigenvector (frame) channel is
linear, its transfer coefficient $2.081$ constant to $0.02\%$ over four and a
half decades. The eigenvalue channel is quadratic, with $k$ between
$4.3\times10^{-12}$ and $6.3\times10^{-11}$. Both coefficients were deposited
before the $2\times2$ run.

The first-order projection is demonstrated directly: the fusion-character
residual is invariant under $300$ random ground-space frame choices (relative
scatter $6.2\times10^{-9}$), because the dominant lattice error --- a
fusion-character deviation of $3.3\times10^{-3}$ --- takes the same value on the
degenerate pair to within $4.8\times10^{-11}$ and so projects out at first
order.

What makes this evidence rather than a fit is that the incomplete version
fails. A pre-registered single-channel formula, keeping only the frame term,
predicted the $2\times2$ residual at $4.4\times10^{-15}$; the measured value is
$6.4\times10^{-10}$ to $1.8\times10^{-9}$, six orders of magnitude larger. Only
the full two-channel law is consistent with both sizes: over all points its
conservative prediction interval, $2.0\times10^{-10}$ to $2.9\times10^{-9}$,
contains the entire measured interval, the worst-case edge ratio being $9.2$.
The dominant channel even switches between sizes --- the frame term on
$3\times3$, from LOBPCG ground states, the quadratic term on $2\times2$, from
dense diagonalization --- and a counterintuitive consequence holds: although
the $2\times2$ loops are $6.8\times$ less ideal, the extraction degrades only
$3.1\times$, because on $2\times2$ the frame is essentially exact. The scale $s$
is set by the deposited loop ideality and is never a free knob.

\subsection{What remains a convention: splits versus labeling}
\label{sec:splits}

One genuine limit survives, and it is a property of the theory, not of the
method. The doubled-Fibonacci theory is achiral ($c=0$), and complex
conjugation $K$ is an exact antiunitary symmetry of the real lattice model
($\overline{T}=P\,T\,P$ for a permutation $P$). The rotation channel therefore
\emph{separates} the chiral pair but does not \emph{label} it: which of $\tau$,
$\bar\tau$ is called which remains a convention, for this method and for any
other, Francuz--Dziarmaga included. Empirically the run lands on either
convention indifferently --- the conjugate match is mixed, $2$ of the $8$
distinct solutions on $3\times3$ and $5$ of $9$ on $2\times2$ --- which is the
prediction, not a contradiction. This $\mathbb{Z}_2$ relabeling is a different
thing from the fusion-phase freedom of Sec.~\ref{sec:unique}: the latter is
trivial for doubled Fibonacci, the former is intrinsic to any achiral theory.

\subsection{What is not used as a certificate}
\label{sec:tube}

We do not use a tube-algebra cross-check~\cite{Bultinck2017} as an independent
certificate, although it would be natural to reach for. For a doubled theory it
is circular: $Z(\mathcal{C})=\mathcal{C}\boxtimes\bar{\mathcal{C}}$ is a
theorem, so any theory-side center computation necessarily factorizes into
$S_{\mathrm{Fib}}\otimes\overline{S_{\mathrm{Fib}}}$ (verified here at
$1.1\times10^{-16}$). ``Two theory methods agree'' is then guaranteed by the
theorem and checks nothing. We cite~\cite{Bultinck2017} as a
methods and context reference only.

\begin{table*}[tp]
\caption{\label{tab:numbers}Load-bearing quantities, with evidence class and
scope. Class (a) is the vacuum row, fixed by two distinct lattice routes; (b)
the loop route and the consistency checks it supports; (c) the ground-state
degeneracy, a separate lock; (d) the point-group rotation channel that resolves
the full signed $S$ and $T$. The (b) rows record what the loop route measures
and, by measuring it, cannot reach; the (d) rows record what resolves it.
Values are reproduced from the deposited artifacts (see Data and Code
Availability).}
\begin{ruledtabular}
\scriptsize
\begin{tabular}{llll}
Quantity & Value & Class & Scope \\
\colrule
Vacuum row $|S_{0a}|=d_a/D$ & $\{0.2764,0.4472,0.4472,0.7236\}$ & (a) two lattice routes & $3\times3$ \\
\quad Verlinde--Casimir vs.\ theory & $1.1\times10^{-16}$ & (a) one route, no analytic $F/R$ & $3\times3$ \\
\quad route to route & $<10^{-9}$ & (a) operator $\wedge$ Verlinde--Casimir & $3\times3$ \\
$S_{00}$ & $0.2763932$, error $1.7\times10^{-15}$; $=1/D$ & (a) lattice & $2\times2$ \\
$\mathrm{GSD}$, gap & $4$, $1.0$ & (c) separate lock & all three sizes \\
$E_0$ (extensive, \emph{not} inv.) & $-4$ ($2\times2$), $-9$ ($3\times3$) & (c) context & size-specific \\
Gauge invariance of $|S|$ & $2.4\times10^{-15}$ ($2\times2$), $5.6\times10^{-15}$ ($2\times3$) & soundness & $24\times O(4)$ \\
$U(4)$ control & fails at $0.603$ & discriminator & --- \\
Size robustness of $|S|$ & $\le2.5\times10^{-11}$ & soundness & $175/2250/106250$ \\
Hexagon / pentagon / Levin--Wen & $2.4\times10^{-16}$ / $0$ / ${\sim}10^{-16}$ & consistency locks & naive: $1.18$ \\
\colrule
Topological spins $T$ & $\{0,0,\pm4\pi/5\}$; Verlinde resid.\ $3.98\times10^{-11}$ & (b) consistency; measured in (d) & lattice $S$ \\
Loop algebra $\langle M_i,M_j\rangle$ & $\dim=10=1^{2}+3^{2}$ & (b) loop route, measured limit & $3\times3$ \\
Antisym.\ split; chiral capacity & $1.2\times10^{-16}$, $7.5\times10^{-17}$; $0.0$ & (b) loop route, measured & $3\times3$ \\
Raw signed $S$, loop route & sym $1.46$; gauge $1.08$; $(ST)^{3}{-}S^{2}\,1.27$; lock \textsc{false} & (b) loop route fails, measured & $3\times3$ \\
Selector-backed old route & $1.16\times10^{-10}$ vs.\ theory & superseded: sign supplied & not lattice-borne \\
\colrule
$C_3$ orthogonality dev. & $2.96\times10^{-15}$, $4.41\times10^{-15}$ & (d) genuine symmetry & $3\times3$/$2\times2$ \\
Algebra jump with $C_3$ & $10\to16$ ($M_R\notin$ loops) & (d) rotation $\ne$ loops & $3\times3$ \\
Signed $S,T$ extraction & $12$ sol.; res.\ $2.9$--$4.8{\times}10^{-10}$; $T{=}\{e^{\mp4\pi i/5},1,1\}$ & (d) rotation channel & $3\times3$ \\
Signed $S,T$ extraction & $9$ sol.; res.\ $6.4{\times}10^{-10}$--$1.8{\times}10^{-9}$ & (d) rotation channel & $2\times2$ \\
$\mathbb{Z}_2$ convention (conj.\ match) & mixed: $2/8$, $5/9$ & (d) intrinsic to achiral theory & $3\times3$/$2\times2$ \\
\colrule
Frame transfer (eigenvector ch.) & $2.081$, linear, $0.02\%$ over $4.5$ dec. & error budget & model \\
$k$ (eigenvalue channel) & $4.3\times10^{-12}$--$6.3\times10^{-11}$, quadratic & error budget & model \\
Loop ideality (input scale) & $2.2\times10^{-2}\to3.3\times10^{-3}$ & deposited, not a knob & $2\times2$/$3\times3$ \\
Two-channel prediction & $2.0{\times}10^{-10}$--$2.9{\times}10^{-9}\supseteq$ meas.\ $6.4{\times}10^{-10}$--$1.8{\times}10^{-9}$ & conservative & both \\
\end{tabular}
\end{ruledtabular}
\end{table*}

\section{A no-go for quasi-one-dimensional anyon exchange}
\label{sec:nogo}

\subsection{Scope}
\label{sec:nogoscope}

We report a bounded negative result. On the one-dimensionally mapped anyonic
ladder of Singh, Pfeifer, Vidal and Brennen~\cite{Singh2014,Ayeni2016}, the
dynamical, Hamiltonian-driven transport of a single mobile $\sigma$ anyon
around a pinned object \emph{cannot} realize the non-Abelian Ising braid
matrix. The closed-loop holonomy is exactly the identity, and the commutator of
two independent loop holonomies vanishes,
\begin{equation}
  [U(C_1),\,U(C_2)] = 0\,.
  \label{eq:nogo}
\end{equation}
We establish this through two independent obstructions, test three escape
routes (all null, with no tuning), and --- the point that makes a null number
trustworthy --- verify with a machinery-alive control that the vanishing is
physics and not a dead numerical pipeline.

This is a statement about \emph{this model}. It is not a claim that non-Abelian
braiding is impossible in general, nor that a two-dimensional geometry cannot
realize it. The static certificate of Sec.~\ref{sec:c1} is untouched by it.

\subsection{The result}
\label{sec:nogoresult}

Let $G$ be the (position $\times$ leg) graph of the mapped ladder, with the
mobile $\sigma$ living on a vertex. The construction has four load-bearing
hypotheses: (i) $G$ is simply connected, so that all cycles are contractible;
(ii) the braid $R$-phase resides purely in the bond amplitudes,
$x=R^{\sigma\sigma}_{\Id}J$ and $y=R^{\sigma\sigma}_{\psi}J$, identically on
both legs; (iii) a single mobile anyon traverses a closed loop; and (iv) the
encircled object is positionless, living in the fusion space rather than on
$G$.

Under these hypotheses the argument is short and exact: total charge is fixed,
so each fusion sector is one-dimensional; the Berry connection is therefore
Abelian $U(1)$ per sector; the net $R$-flux vanishes on every contractible
cycle; the $R$-phase is pure gauge on $G$; hence $U(C)=\Id$ exactly, and
Eq.~\eqref{eq:nogo} follows.

Numerically: plaquette $R$-flux $=0$; gauge-reconstruction residual
$1.7\times10^{-16}$; spectrum invariance under the gauge transform
$2.1\times10^{-15}$; the full $H_1$ cycle-basis maximal flux $=0$ on
$N=3,5,7$; and a dense exact-diagonalization adiabatic transport giving
Eq.~\eqref{eq:nogo} with operator and Frobenius norm $0$ and
$U(C_1)=U(C_2)=\Id$.

The same statement is deposited independently at the \emph{structural} level:
building each holonomy directly from the exact $R$-symbols as
$U(C)=\mathrm{diag}(e^{i\,\mathrm{net}_{\Id}},e^{i\,\mathrm{net}_{\psi}})$
gives $\Fnorm{U(C_1)-\Id} = \Fnorm{U(C_2)-\Id} = 0$ and
$\Fnorm{[U(C_1),U(C_2)]} = \Opnorm{[U(C_1),U(C_2)]} = 0$. We label the two
levels separately and do not merge them into one ``verified'' number: the
deposited structural route and the dense-ED adiabatic transport are two levels
of evidence for one statement, the latter being the independent confirmation of
the former.

Two positive controls keep the null honest. The holonomy builder is live: a
loop with \emph{nonzero} net winding returns $\Fnorm{U_W-\Id}=2\ne0$. The
commutator builder is live: the same norm code path, fed an $F$-dressed
off-diagonal pair, returns $\sqrt2\ne0$, while the $F=\Id$ anti-fudge collapses
it to exactly $0$. We note explicitly what we did \emph{not} use as a control:
two winding loops over the same bond-phase pair are powers of one Abelian
generator, so their commutator vanishes trivially and would prove nothing. An
early version of the control was exactly that vacuous, and was discarded.

\subsection{Two independent obstructions}
\label{sec:obstructions}

\emph{Topological.} A one-dimensional strip cannot enclose a point defect: any
closed walk crosses each bond a net zero number of times, so the $R$-phase
cancels and is gauge removable. Here
$\mathrm{net}_{\Id}=\mathrm{net}_{\psi}=0.0$. This is independently consistent
with the graph-braid-group classification of anyons on quantum wire
networks~\cite{ConlonSlingerland2022,Maciazek2023}: those results indicate that
a network which is not triconnected does not support planar encircling braiding.

\emph{Algebraic (superselection).} The total Ising fusion charge
$\{\Id,\psi\}$ is an exact superselection sector: both hopping and pinning have
zero matrix element between the $\Id$ and $\psi$ blocks. The non-Abelian
generator $2\sigma_y$ requires an $\Id\!\leftrightarrow\!\psi$ rotation --- the
$F$-move --- and no local or spatial operation generates it. Over $10^{4}$
random transports diagonal in $\{\Id,\psi\}$ the commutator vanishes: the
deposited value is exactly $0$, as it must be, since operators diagonal in a
common basis commute by construction. An independent reimplementation of the
\emph{same} test reproduces this to its floating-point floor
($1.57\times10^{-16}$, a one-ulp artifact of the complex matrix multiply, not a
physical residue, and not a second independent check).

The \emph{mechanism} is confirmed by a separate discriminator: the non-Abelian
commutator exists only through the $F$-move. With the true (Hadamard) $F$ one
finds $\Fnorm{[s_1,s_2]}=\sqrt2$; setting $F=\Id$ forces it to exactly $0$.

The two obstructions are the same cause seen two ways, one topological and one
algebraic, and so constitute two independent proof paths for a single no-go.
The step ``flux $=0 \Rightarrow$ trivial holonomy'' holds only for Abelian
connections; it is airtight here because each fixed-charge sector is scalar per
channel.

\subsection{Escape routes: live-tested, anti-fudge-gated, all null}
\label{sec:escapes}

\begin{table}[tbp]
\caption{\label{tab:escapes}The three escape routes, each relaxing a
load-bearing hypothesis of Sec.~\ref{sec:nogoresult}. All are null; none
required tuning or a convention flip.}
\begin{ruledtabular}
\scriptsize
\begin{tabular}{lll}
Escape & Relaxed & Outcome \\
\colrule
Two-$\sigma$ exchange & (iii), (iv) & Exchange is diagonal/Abelian; \\
 & & $[B_{12},B_{23}]=0$ in both norms. \\
 & & Cause: superselection. \\
Periodic ring & (i) & A winding phase survives but is \\
 & & unchanged under bond shift: an \\
 & & Aharonov--Bohm hole flux, not \\
 & & defect encircling; and Abelian. \\
Direct $B$ operator & bypasses space & Oracle artifact: \\
 & & $\Fnorm{[B_{1m},B_{2m}]}=\sqrt2$ \\
 & & for \emph{every} identification, \\
 & & i.e.\ zero emergent content. \\
\end{tabular}
\end{ruledtabular}
\end{table}

Table~\ref{tab:escapes} summarizes. The first relaxes hypotheses (iii) and
(iv): an $N=4$ dense-ED two-$\sigma$ exchange, on blocks of dimension $48$
each, with the graph fully connected so that a path exists. The exchange is
diagonal and Abelian ($\varphi_\psi-\varphi_{\Id}\approx\pi$) and the braid
generators commute in both norms; the cause is again superselection. The second
relaxes (i): on an $L=6$ ring a winding phase does survive and is not gauge
removable, but shifting the bond leaves it unchanged, which identifies it as an
Aharonov--Bohm flux through the lattice hole rather than defect encircling ---
and it is Abelian in any case. The third bypasses space entirely, building $B$
from $F$ and $R$ on the fusion space; it returns the same nonzero value for
\emph{every} identification, which marks it as an oracle rather than a
demonstration, and the corresponding lock correctly labels it as such.

Each route carried its own controls: machinery-alive, gauge invariance, and an
oracle-versus-simulation separation for the first; a braid-free null and an
injected $\pi/2$ read back for the second; and the oracle lock for the third.

\subsection{Machinery-alive control}
\label{sec:machinery}

The identical singular-value/Wilczek--Zee transport
pipeline~\cite{WilczekZee1984}, applied to a synthetic control carrying a
genuine $\pi$ rotation (a $2\pi$-solid-angle loop), returns
$U_{\mathrm{control}}=-\Id$, with
\begin{equation}
  \Fnorm{U_{\mathrm{control}}-\Id} = 2\sqrt2 \approx 2.828427\,.
  \label{eq:alive}
\end{equation}
A dead or short-circuited pipeline could not produce this known nontrivial
answer. The vanishing holonomy of the physical transport is therefore a genuine
property of the model, not an artifact of the numerics.

\subsection{Oracle versus lattice: the correct labeling}
\label{sec:oracle}

The abstract Ising braid \emph{representation} does contain the non-Abelian
structure: the single-exchange commutator $[s_1,s_2]$ has Frobenius norm
$\sqrt2=\Fnorm{\sigma_y}$, and the monodromy $[s_1^{2},s_2^{2}]$ has
Frobenius norm $2\sqrt2=\Fnorm{2\sigma_y}$. These are properties of the
representation --- the oracle --- obtained by substituting the known $F$ and
$R$ data. They answer a \emph{different question} from the dynamical lattice
transport, which returns $0$. Both numbers are correct, and they must be
labeled separately, or a spurious contradiction is invited.

The same discipline applies to the energetics. The oracle's two-island Ising
chain has $E_{\mathrm{GS}}=-(1+\sqrt2)$ and gap $\sqrt2$, with the $N=3$
spectrum $\{\pm\cot(\pi/8),\pm1,\pm\tan(\pi/8)\}$ reproduced at $10^{-15}$;
this is a different object from the doubled-Fibonacci lattice of
Sec.~\ref{sec:c1}, whose ground-state energy is extensive ($E_0=-4$ on
$2\times2$, $-9$ on $3\times3$). We never mix them.

\subsection{What this section claims, and does not}
\label{sec:nogoclaims}

We claim that on this specific one-dimensionally mapped ladder, the spatial
transport of a single anyon provably cannot realize the non-Abelian braid
matrix, and that the mechanism --- gauge removability together with
$\{\Id,\psi\}$ superselection --- is understood.

We do not claim that non-Abelian braiding is impossible in general, nor that a
two-dimensional geometry cannot realize it. And we state plainly that
$M_{\tau\tau}=-1/\varphi^{2}$ of Eq.~\eqref{eq:target} was the \emph{aim} of
the dynamic half and was never measured; this section is the reason why. It
must not be confused with the Wilson-loop eigenvalue $+1/\varphi^{2}$ appearing
in the static construction of Sec.~\ref{sec:c1}, which is a different quantity
with the opposite sign.

\section{Discussion and outlook}
\label{sec:discussion}

\subsection{The compass}
\label{sec:compass}

Read together, the two halves point in one direction. Spatial braiding does
not live in one-dimensional single-particle transport. That is a transferable
lesson rather than a local accident, because both obstructions of
Sec.~\ref{sec:obstructions} are structural.

The same structural reading explains the static half. The lattice carries the
real, basis-invariant modular content because that content is encoded in the
operator algebra of the string-net; it carries the chiral sign as well, but
through the point-group rotation, not through the Wilson loops, whose
real-symmetric holonomies are phase-free. The cut of Sec.~\ref{sec:cut} is not
where the lattice stops but where the loop route does, and the rotation channel
is what crosses it. More broadly, the channel suggests reading modular data off
a finite system from a symmetry of the system at rest, rather than by
transporting a particle around it --- a perspective we note without claiming it
as a demonstrated application.

\subsection{Relation to prior work}
\label{sec:priorart}

We make no priority claim, and one piece of prior work sits close enough that
we state it precisely rather than let a reader discover it. Francuz and
Dziarmaga~\cite{FrancuzDziarmaga2020} extract the topological $S$ and $T$
matrices of non-Abelian string-net models --- \emph{including the Fibonacci
model treated here} --- from an infinite projected entangled pair state, using
matrix-product-operator symmetries and their fusion rules to build anyon-flux
projectors. The modular data of the Fibonacci string-net are therefore not new
with this work, and we do not claim them.

What differs is the method, and with it the question that can be asked. Theirs
is a tensor-network computation with exact renormalization-group fixed-point
tensors, evaluated in the infinite-cylinder limit; ours is a classically exact
diagonalization of a finite lattice, in which a Wilson string is genuinely
moved outside the ground space. That difference is not
cosmetic here: it is what makes the obstruction of Sec.~\ref{sec:wedderburn}
\emph{measurable} --- one can ask what a finite real-symmetric holonomy does and
does not resolve, and get a number --- and it is what makes the transport
question of Sec.~\ref{sec:nogo} well posed at all. Neither question is posed by
an infinite-cylinder ansatz.

Earlier work in the same family determined topological order for the Abelian
case~\cite{Francuz2020} and characterized matrix-product-operator symmetries
variationally~\cite{Francuz2021}; generalized string-net degeneracies have been
analyzed in detail~\cite{RitzZwilling2024}; modular matrices as topological
order parameters~\cite{HeMoradiWen2014} and their extraction by shearing
lattices~\cite{YouCheng2015} are established. The rotation channel we use is the
construction of Li and Mong~\cite{LiMong2022}, whose theorem also bounds what
the loop overlaps alone can give; exact diagonalization of a Fibonacci
string-net is itself not new, having been used to study the golden-chain
transition~\cite{Schulz2013}, and minimal quantum circuits for Fibonacci anyons
have recently been constructed~\cite{Bseiso2024}. Our contribution is the
finite-lattice demonstration of the rotation channel, the quantified
two-channel error budget, and the measured loop-algebra bound --- not the
modular data, and not the channel itself.

On the hardware side, Fibonacci statistics have been realized as digital gate
sequences~\cite{Xu2024}, and non-Abelian braiding of graph vertices and
non-Abelian topological order have been demonstrated on superconducting and
trapped-ion processors respectively~\cite{Andersen2023,Iqbal2024}, and
Minev \emph{et al.}~\cite{Minev2025} prepare a Fibonacci string-net condensate
on a superconducting processor and braid its anyons there. These are published
positive results and we cite them as such, not as a backdrop of absence. The
line to the present work runs the opposite way from the one a reader might
expect, so we draw it explicitly: their question is the \emph{physical
realization} of the condensate --- which Ref.~\cite{Minev2025} opens by
describing as having remained elusive, and which that work then addresses ---
whereas ours is a classically exact computation on a finite lattice. The two are
complementary, and the present paper contributes nothing to resolving the
realization question.

To our knowledge the specific combination reported here --- a real-space,
out-of-ground-state, chirality-discriminating string in a doubled-Fibonacci
Levin--Wen model, the \emph{measured} obstruction of
Sec.~\ref{sec:wedderburn}, and the paired no-go of Sec.~\ref{sec:nogo} --- is
independent of that work, but we have not crawled the prior art exhaustively,
and we make no claim of being first at anything. Our no-go is independently
consistent with the wire-network
classification~\cite{ConlonSlingerland2022,Maciazek2023}, which reaches a
compatible conclusion from a different direction. The universality-theoretic
backdrop for Fibonacci anyons is due to Freedman, Larsen and
Wang~\cite{FLW2002}, with density results due to Kuperberg~\cite{Kuperberg2009}.

\subsection{The remaining niche}
\label{sec:niche}

One direction from the original program survives, and we place it here rather
than in the body because it is off-model. A classically exact, static
Hamiltonian, velocity-resolved diabatic-plus-leakage budget for a universal
(Fibonacci) anyon would be a genuine target. It belongs to a separate, future
strand and not to this paper. We flag it as an open direction, not as a
result.

\section*{Data and Code Availability}
\label{sec:data}

The raw JSON outputs and the reproduction script supporting the values
reported in this paper are deposited at
\href{https://doi.org/10.5281/zenodo.21362245}{doi:10.5281/zenodo.21362245}.
Paper, figures, and data are released under the Creative Commons Attribution
4.0 International (CC~BY~4.0) license; the deposited source code is released
under the Apache License, Version 2.0
(see \texttt{LICENSE} and \texttt{LICENSE-CODE}).

The deposit contains the load-bearing quantities of Table~\ref{tab:numbers}
together with their provenance. For the static certificate and the no-go:
\texttt{p6\_nogo.json} (the structural no-go route, its two positive controls, and
the machinery-alive control), \texttt{p6\_manifold.json} (the $N=3$ dense-ED
spectrum), \texttt{p6\_oracle.json} (the algebra oracle, explicitly labeled as such
and not as a measurement), \texttt{p6\_c1\_modular.json} (the class-(a)--(c)
quantities and the measured loop-route limit, each carrying its evidence class
and per-entry provenance), and \texttt{p6\_modular\_reference.py}, a
self-contained script (\textsc{numpy} only, no engine import, no stored oracle)
that rebuilds the analytic doubled-Fibonacci modular data from scratch and
checks the analytic targets, including the Gauss sum $p_{+}=D$ and the rejection
of the order-3 alternative at $5.236$. For the rotation channel:
\texttt{p6\_c3\_results.json} (the $C_3$ symmetry and the algebra jump $10\to16$),
\texttt{p6\_lattice\_extraction.json} and \texttt{p6\_extraction\_2x2.json} (the signed
$S$ and $T$ extractions on the two sizes), \texttt{p6\_reconciliation\_test.json}
and \texttt{p6\_loop\_ideality.json} (the two-channel error budget and the input
scale), together with the rotation-channel extractor \texttt{p6\_extract\_c3.py}
that produced them, which takes only measured matrices as input, and its
pre-registered referee harness \texttt{p6\_verify\_patch.py} (the five acceptance
gates and the scaling gate). The error budget is reproducible from this record as
well: \texttt{p6\_reconciliation\_test.py} is the deterministic generator of
\texttt{p6\_reconciliation\_test.json} and takes only the deposited loop ideality
and \texttt{p6\_extract\_c3.py} as input. On both sizes the measured input matrices are
deposited, each with its $C_3$ rotation matrix for the center $(0,0)$:
\texttt{p6\_V\_gs}, \texttt{p6\_Mgs\_i}, \texttt{p6\_Mgs\_j} and \texttt{p6\_M\_R} in
\texttt{3x3} and \texttt{2x2} versions, the latter set with the provenance
record \texttt{p6\_manifest\_2x2.json}. The extraction therefore runs from this
record: applying \texttt{p6\_extract\_c3.py} to the deposited Wilson and rotation
matrices, at the deposited loop ideality as input scale and with 240 restarts,
returns the four solutions of center $(0,0)$ recorded in
\texttt{p6\_lattice\_extraction.json} and the three recorded in
\texttt{p6\_extraction\_2x2.json}, reproducing their residuals and sign structure to
all printed digits; the README gives both calls verbatim. Both input scales are
deposited, not fitted. For $3\times3$ we additionally deposit one rotation
matrix per permutation class, \texttt{p6\_M\_R\_3x3\_center20.npy} (center
$(2,0)$, the second matrix behind the recorded run) and
\texttt{p6\_M\_R\_3x3\_center10.npy} (center $(1,0)$, the class the run did not
visit), with \texttt{p6\_center\_independence.py}, which checks from the record
that the projected rotation does not depend on the center and that the extractor
returns the same block structure and the recorded residuals on each; their
construction from the string-net basis is not part of this record. The further
$2\times2$ centers use matrices built from the lattice basis and are not
deposited. Note that the
Wilson matrix carrying the extraction is \texttt{p6\_Mgs\_j\_3x3.npy} on one size
and \texttt{p6\_Mgs\_i\_2x2.npy} on the other: the convention does not carry over,
and the README states why. The residuals
are not an output of the extractor. They are formed in a separate scoring step
against the analytic doubled-Fibonacci target that
\texttt{p6\_modular\_reference.py} rebuilds from scratch; the driver performing
that step is released with this record, while the generation of the measured
matrices themselves is not, since the exact-diagonalization engine that produces
the ground-state and loop matrices is a separate, heavier computation. A minimal
ED path (\texttt{p6\_ed\_path.py}) and its geometry record
(\texttt{p6\_geometry\_2x2.json}) identify the row ordering of the deposited
basis; that path is not the production engine.

\section*{Use of AI tools}

In the research, processing, and writing of this paper and its results I worked together with generative AI tools, in practice a system of multiple coordinated AI instances that I set up and orchestrate (large language models, mainly Claude, by Anthropic, inside Claude Code). At their current context sizes I found it far more effective to work with several specialized instances, each with its own role and its own harness of rules and parameters that I designed and refined through feedback, than to load a single instance with all of the material; for my workflow that would have been inefficient, though this depends on the individual implementation. I lead this collaboration: I choose the research directions, set the goals, and make the final decisions in open exchange with the AI, learning actively as the work proceeds. The AI carries out the drafting, including the mathematical and technical parts, the numerical computation, and the literature search, under my direction. The AI works autonomously only task by task, within the structure I develop through feedback: it completes a task, and at open questions that need me it stops until the point is settled before the next step. Along the way I witness and take many of the decisions that shape the path, and it is common for me to spot things that need improvement. The work spans many separate runs, and a single simulation or build task alone can take up to an hour, so it could not happen all together in one autonomous run; and had I let the AI do all of it together alone, even if it is possible, it would no longer be my work but the AI's.

I run multiple verifications at the different stages of the work and one before release, including cross-checks with an unrelated AI model from a different company, and all references are checked against the original sources. In the end what matters are human eyes, a principle that is itself written into the parameters of my system: I reach out to experts after publishing for review and feedback, so I learn what is solid and what must be corrected or falsified. My scripts for reproduction and review are released with this record. These tools are not authors; I am the author, and I take full responsibility for all scientific content and decisions leading to these results and their publication.

\section*{Acknowledgments}
No external funding was received. The author declares no competing
interests.

\begin{thebibliography}{99}

\bibitem{Xu2024}
S.~Xu, Z.-Z.~Sun, K.~Wang, H.~Li, Z.~Zhu, H.~Dong, J.~Deng, X.~Zhang, J.~Chen,
Y.~Wu, C.~Zhang, F.~Jin, X.~Zhu, Y.~Gao, A.~Zhang, N.~Wang, Y.~Zou, Z.~Tan,
F.~Shen, J.~Zhong, Z.~Bao, W.~Li, W.~Jiang, L.-W.~Yu, Z.~Song, P.~Zhang,
L.~Xiang, Q.~Guo, Z.~Wang, C.~Song, H.~Wang, and D.-L.~Deng,
\emph{Non-Abelian braiding of Fibonacci anyons with a superconducting
processor},
Nat.\ Phys.\ \textbf{20}, 1469 (2024);
\href{https://arxiv.org/abs/2404.00091}{arXiv:2404.00091}.

\bibitem{LevinWen2005}
M.~A.~Levin and X.-G.~Wen,
\emph{String-net condensation: A physical mechanism for topological phases},
Phys.\ Rev.\ B \textbf{71}, 045110 (2005);
\href{https://arxiv.org/abs/cond-mat/0404617}{arXiv:cond-mat/0404617}.

\bibitem{LiMong2022}
Z.~Li and R.~S.~K.~Mong,
\emph{Detecting topological order from modular transformations of ground states
on the torus},
Phys.\ Rev.\ B \textbf{106}, 235115 (2022);
\href{https://arxiv.org/abs/2203.04329}{arXiv:2203.04329}.

\bibitem{Schulz2013}
M.~D.~Schulz, S.~Dusuel, K.~P.~Schmidt, and J.~Vidal,
\emph{Topological phase transitions in the golden string-net model},
Phys.\ Rev.\ Lett.\ \textbf{110}, 147203 (2013);
\href{https://arxiv.org/abs/1212.4109}{arXiv:1212.4109}.

\bibitem{Bseiso2024}
S.~Bseiso, J.~Pommerening, R.~R.~Allen, S.~H.~Simon, and L.~Hormozi,
\emph{Minimal quantum circuits for simulating Fibonacci anyons},
\href{https://arxiv.org/abs/2407.21761}{arXiv:2407.21761}.

\bibitem{Singh2014}
S.~Singh, R.~N.~C.~Pfeifer, G.~Vidal, and G.~K.~Brennen,
\emph{Matrix product states for anyonic systems and efficient simulation of
dynamics},
Phys.\ Rev.\ B \textbf{89}, 075112 (2014);
\href{https://arxiv.org/abs/1311.0967}{arXiv:1311.0967}.

\bibitem{Ayeni2016}
B.~M.~Ayeni, S.~Singh, R.~N.~C.~Pfeifer, and G.~K.~Brennen,
\emph{Simulation of braiding anyons using matrix product states},
Phys.\ Rev.\ B \textbf{93}, 165128 (2016);
\href{https://arxiv.org/abs/1509.00903}{arXiv:1509.00903}.

\bibitem{Verlinde1988}
E.~Verlinde,
\emph{Fusion rules and modular transformations in 2D conformal field theory},
Nucl.\ Phys.\ B \textbf{300}, 360 (1988).

\bibitem{Vafa1988}
C.~Vafa,
\emph{Toward classification of conformal theories},
Phys.\ Lett.\ B \textbf{206}, 421 (1988).

\bibitem{RSW2009}
E.~Rowell, R.~Stong, and Z.~Wang,
\emph{On classification of modular tensor categories},
Commun.\ Math.\ Phys.\ \textbf{292}, 343 (2009);
\href{https://arxiv.org/abs/0712.1377}{arXiv:0712.1377}.

\bibitem{Bantay1997}
P.~Bantay,
\emph{The Frobenius-Schur indicator in conformal field theory},
Phys.\ Lett.\ B \textbf{394}, 87 (1997);
\href{https://arxiv.org/abs/hep-th/9610192}{arXiv:hep-th/9610192}.

\bibitem{Kitaev2006}
A.~Kitaev,
\emph{Anyons in an exactly solved model and beyond},
Ann.\ Phys.\ \textbf{321}, 2 (2006);
\href{https://arxiv.org/abs/cond-mat/0506438}{arXiv:cond-mat/0506438}.

\bibitem{NgSchauenburg2010}
S.-H.~Ng and P.~Schauenburg,
\emph{Congruence subgroups and generalized Frobenius-Schur indicators},
Commun.\ Math.\ Phys.\ \textbf{300}, 1 (2010);
\href{https://arxiv.org/abs/0806.2493}{arXiv:0806.2493}.

\bibitem{DongLinNg2015}
C.~Dong, X.~Lin, and S.-H.~Ng,
\emph{Congruence property in conformal field theory},
Algebra Number Theory \textbf{9}, 2121 (2015);
\href{https://arxiv.org/abs/1201.6644}{arXiv:1201.6644}.

\bibitem{Freedman2012}
M.~H.~Freedman, J.~Gukelberger, M.~B.~Hastings, S.~Trebst,
M.~Troyer, and Z.~Wang,
\emph{Galois conjugates of topological phases},
Phys.\ Rev.\ B \textbf{85}, 045414 (2012);
\href{https://arxiv.org/abs/1106.3267}{arXiv:1106.3267}.

\bibitem{Bultinck2017}
N.~Bultinck, M.~Mari\"en, D.~J.~Williamson, M.~B.~\c{S}ahino\u{g}lu,
J.~Haegeman, and F.~Verstraete,
\emph{Anyons and matrix product operator algebras},
Ann.\ Phys.\ \textbf{378}, 183 (2017);
\href{https://arxiv.org/abs/1511.08090}{arXiv:1511.08090}.

\bibitem{ConlonSlingerland2022}
M.~Conlon and J.~K.~Slingerland,
\emph{Compatibility of braiding and fusion on wire networks},
Phys.\ Rev.\ B \textbf{108}, 035150 (2023);
\href{https://arxiv.org/abs/2202.08207}{arXiv:2202.08207}.

\bibitem{Maciazek2023}
T.~Maciazek, M.~Conlon, G.~Vercleyen, and J.~K.~Slingerland,
\emph{Extending the planar theory of anyons to quantum wire networks},
SciPost Phys.\ \textbf{18}, 074 (2025);
\href{https://arxiv.org/abs/2301.06590}{arXiv:2301.06590}.

\bibitem{WilczekZee1984}
F.~Wilczek and A.~Zee,
\emph{Appearance of gauge structure in simple dynamical systems},
Phys.\ Rev.\ Lett.\ \textbf{52}, 2111 (1984).

\bibitem{Francuz2020}
A.~Francuz, J.~Dziarmaga, G.~Vidal, and L.~Cincio,
\emph{Determining topological order from infinite projected entangled pair
states},
Phys.\ Rev.\ B \textbf{101}, 041108 (2020);
\href{https://arxiv.org/abs/1910.09661}{arXiv:1910.09661}.

\bibitem{FrancuzDziarmaga2020}
A.~Francuz and J.~Dziarmaga,
\emph{Determining non-Abelian topological order from infinite projected
entangled pair states},
Phys.\ Rev.\ B \textbf{102}, 235112 (2020);
\href{https://arxiv.org/abs/2008.06391}{arXiv:2008.06391}.

\bibitem{Francuz2021}
A.~Francuz, L.~Lootens, F.~Verstraete, and J.~Dziarmaga,
\emph{Variational methods for characterizing matrix product operator
symmetries},
Phys.\ Rev.\ B \textbf{104}, 195152 (2021);
\href{https://arxiv.org/abs/2107.05265}{arXiv:2107.05265}.

\bibitem{Minev2025}
Z.~K.~Minev, K.~Najafi, S.~Majumder, J.~Wang, A.~Stern, E.-A.~Kim, C.-M.~Jian,
and G.~Zhu,
\emph{Realizing string-net condensation: Fibonacci anyon braiding for universal
gates and sampling chromatic polynomials},
Nat.\ Commun.\ \textbf{16}, 6225 (2025);
\href{https://arxiv.org/abs/2406.12820}{arXiv:2406.12820}.

\bibitem{RitzZwilling2024}
A.~Ritz-Zwilling, J.-N.~Fuchs, S.~H.~Simon, and J.~Vidal,
\emph{Topological and nontopological degeneracies in generalized string-net
models},
Phys.\ Rev.\ B \textbf{109}, 045130 (2024);
\href{https://arxiv.org/abs/2309.00343}{arXiv:2309.00343}.

\bibitem{HeMoradiWen2014}
H.~He, H.~Moradi, and X.-G.~Wen,
\emph{Modular matrices as topological order parameter by gauge symmetry
preserved tensor renormalization approach},
Phys.\ Rev.\ B \textbf{90}, 205114 (2014);
\href{https://arxiv.org/abs/1401.5557}{arXiv:1401.5557}.

\bibitem{YouCheng2015}
Y.-Z.~You and M.~Cheng,
\emph{Measuring modular matrices by shearing lattices},
\href{https://arxiv.org/abs/1502.03192}{arXiv:1502.03192} (unpublished).

\bibitem{Andersen2023}
T.~I.~Andersen \emph{et al.} (Google Quantum AI and Collaborators),
\emph{Non-Abelian braiding of graph vertices in a superconducting processor},
Nature \textbf{618}, 264 (2023);
\href{https://arxiv.org/abs/2210.10255}{arXiv:2210.10255}.

\bibitem{Iqbal2024}
M.~Iqbal, N.~Tantivasadakarn, R.~Verresen, S.~L.~Campbell, J.~M.~Dreiling,
C.~Figgatt, J.~P.~Gaebler, J.~Johansen, M.~Mills, S.~A.~Moses, J.~M.~Pino,
A.~Ransford, M.~Rowe, P.~Siegfried, R.~P.~Stutz, M.~Foss-Feig,
A.~Vishwanath, and H.~Dreyer,
\emph{Non-Abelian topological order and anyons on a trapped-ion processor},
Nature \textbf{626}, 505 (2024);
\href{https://arxiv.org/abs/2305.03766}{arXiv:2305.03766}.

\bibitem{FLW2002}
M.~H.~Freedman, M.~J.~Larsen, and Z.~Wang,
\emph{The two-eigenvalue problem and density of Jones representation of braid
groups},
Commun.\ Math.\ Phys.\ \textbf{228}, 177 (2002);
\href{https://arxiv.org/abs/math/0103200}{arXiv:math/0103200}.

\bibitem{Kuperberg2009}
G.~Kuperberg,
\emph{Denseness and Zariski denseness of Jones braid representations},
Geom.\ Topol.\ \textbf{15}, 11 (2011);
\href{https://arxiv.org/abs/0909.1881}{arXiv:0909.1881}.

\end{thebibliography}

\end{document}
